% GENERATED FILE. Edit canonical lessons, then run npm run generate:books.
% canonical-source-hash: fnv1a64:c595dff20d7adcb3
% canonical-lessons: GU-C251-arth, GU-C251-anuvad, GU-C251-shabdkosh, GU-C251-dhyan-apvu, GU-C251-yad-karvu

\chapter{Words and their meanings}
\label{ch:gu-words-and-their-meanings}
\hlchaptermodality{251}

\begin{chapteropening}
\textbf{By the end of this chapter:} \emph{I can ask about meanings, translations and dictionaries.}

Complete the last lesson of chapter 251: I can ask about meanings, translations and dictionaries.
\end{chapteropening}

\section[\texorpdfstring{\gu{અર્થ} (arth)}{અર્થ (arth)}]{\gu{અર્થ} (arth) — meaning}

\label{lesson:GU-C251-arth}



\begin{quote}
Before the new one: say the Gujarati for to keep safe, then the Gujarati for to print.
\end{quote}

\subsection*{You'll want to know: \gu{અર્થ}}
\textbf{\gu{અર્થ}} — \emph{arth} — \textquotedblleft{}meaning\textquotedblright{}.

\begin{quote}
\textbf{\gu{આ} \gu{શબ્દનો} \gu{અર્થ} \gu{શું} \gu{છે}?} — \emph{ā śabdano arth śũ che?} — \textquotedblleft{}What is the meaning of this word?\textquotedblright{}
\end{quote}

\begin{cousinweb}
\textbf{\gu{અર્થ}} (\emph{arth}) is Sanskrit \emph{ártha}, \textquotedblleft{}aim, thing\textquotedblright{} and later \textquotedblleft{}meaning\textquotedblright{}, taken over whole as a tatsama (CDIAL 638). The same entry lists an inherited Gujarati \emph{āth}, \textquotedblleft{}store of money\textquotedblright{}.
\end{cousinweb}

\begin{grammarlens}[title={What you've built}]
One more word for everyday life at A2.
\end{grammarlens}

\subsection*{Your turn}
\begin{itemize}
\raggedright
  \item \emph{Say it:} \emph{arth}
  \item \emph{Say it:} \emph{arth}, once more
  \item \emph{Say it:} \emph{ā śabdano arth śũ che?} — What is the meaning of this word?
  \item \emph{Recall:} say the Gujarati for to keep safe, then the Gujarati for to print, then say \emph{arth} again
\end{itemize}

\begin{tcolorbox}[breakable,colback=teal!4,colframe=teal!35!black,title={Before you move on}]
What does \gu{અર્થ} mean? (\textquotedblleft{}Meaning\textquotedblright{}.) Say it once more.
\end{tcolorbox}
\section[\texorpdfstring{\gu{અનુવાદ} (anuvād)}{અનુવાદ (anuvād)}]{\gu{અનુવાદ} (anuvād) — translation}

\label{lesson:GU-C251-anuvad}



\begin{quote}
Before the new one: say the Gujarati for to print, then the Gujarati for meaning.
\end{quote}

\subsection*{You'll want to know: \gu{અનુવાદ}}
\textbf{\gu{અનુવાદ}} — \emph{anuvād} — \textquotedblleft{}translation\textquotedblright{}.

\begin{quote}
\textbf{\gu{તે} \gu{પત્રનો} \gu{અનુવાદ} \gu{કરે} \gu{છે}.} — \emph{te patrano anuvād kare che} — \textquotedblleft{}She translates the letter.\textquotedblright{}
\end{quote}

\begin{cousinweb}
\textbf{\gu{અનુવાદ}} (\emph{anuvād}) is Sanskrit \emph{anuvāda}, built from \emph{anu}, \textquotedblleft{}after\textquotedblright{}, and \emph{vāda}, \textquotedblleft{}speaking\textquotedblright{}, taken over whole as a tatsama. CDIAL has no entry for it.
\end{cousinweb}

\begin{grammarlens}[title={What you've built}]
One more word for everyday life at A2.
\end{grammarlens}

\subsection*{Your turn}
\begin{itemize}
\raggedright
  \item \emph{Say it:} \emph{anuvād}
  \item \emph{Say it:} \emph{anuvād}, once more
  \item \emph{Say it:} \emph{te patrano anuvād kare che} — She translates the letter.
  \item \emph{Recall:} say the Gujarati for to print, then the Gujarati for meaning, then say \emph{anuvād} again
\end{itemize}

\begin{tcolorbox}[breakable,colback=teal!4,colframe=teal!35!black,title={Before you move on}]
What does \gu{અનુવાદ} mean? (\textquotedblleft{}Translation\textquotedblright{}.) Say it once more.
\end{tcolorbox}
\section[\texorpdfstring{\gu{શબ્દકોશ} (śabdkoś)}{શબ્દકોશ (śabdkoś)}]{\gu{શબ્દકોશ} (śabdkoś) — dictionary}

\label{lesson:GU-C251-shabdkosh}



\begin{quote}
Before the new one: say the Gujarati for meaning, then the Gujarati for translation.
\end{quote}

\subsection*{You'll want to know: \gu{શબ્દકોશ}}
\textbf{\gu{શબ્દકોશ}} — \emph{śabdkoś} — \textquotedblleft{}dictionary\textquotedblright{}.

\begin{quote}
\textbf{\gu{હું} \gu{શબ્દકોશમાં} \gu{અર્થ} \gu{શોધું} \gu{છું}.} — \emph{hũ śabdkośmā̃ arth śodhũ chũ} — \textquotedblleft{}I look up the meaning in the dictionary.\textquotedblright{}
\end{quote}

A masculine noun.

\begin{cousinweb}
\textbf{\gu{શબ્દકોશ}} (\emph{śabdkoś}), a store of words, is a learned compound of two tatsamas: Sanskrit \emph{śabda}, \textquotedblleft{}sound, word\textquotedblright{} (CDIAL 12298), and \emph{kōśa}, \textquotedblleft{}case, storeroom\textquotedblright{} (CDIAL 3539).
\end{cousinweb}

\begin{grammarlens}[title={What you've built}]
One more word for everyday life at A2.
\end{grammarlens}

\subsection*{Your turn}
\begin{itemize}
\raggedright
  \item \emph{Say it:} \emph{śabdkoś}
  \item \emph{Say it:} \emph{śabdkoś}, once more
  \item \emph{Say it:} \emph{hũ śabdkośmā̃ arth śodhũ chũ} — I look up the meaning in the dictionary.
  \item \emph{Recall:} say the Gujarati for meaning, then the Gujarati for translation, then say \emph{śabdkoś} again
\end{itemize}

\begin{tcolorbox}[breakable,colback=teal!4,colframe=teal!35!black,title={Before you move on}]
What does \gu{શબ્દકોશ} mean? (\textquotedblleft{}Dictionary\textquotedblright{}.) Say it once more.
\end{tcolorbox}
\section[\texorpdfstring{\gu{ધ્યાન} \gu{આપવું} (dhyān āpvũ)}{ધ્યાન આપવું (dhyān āpvũ)}]{\gu{ધ્યાન} \gu{આપવું} (dhyān āpvũ) — to pay attention}

\label{lesson:GU-C251-dhyan-apvu}



\begin{quote}
Before the new one: say the Gujarati for translation, then the Gujarati for dictionary.
\end{quote}

\subsection*{You'll want to know: \gu{ધ્યાન} \gu{આપવું}}
\textbf{\gu{ધ્યાન} \gu{આપવું}} — \emph{dhyān āpvũ} — \textquotedblleft{}to pay attention\textquotedblright{}.

\begin{quote}
\textbf{\gu{વિદ્યાર્થી} \gu{વર્ગમાં} \gu{ધ્યાન} \gu{આપે} \gu{છે}.} — \emph{vidyārthī vargmā̃ dhyān āpe che} — \textquotedblleft{}The students pay attention in class.\textquotedblright{}
\end{quote}

\begin{cousinweb}
\textbf{\gu{ધ્યાન} \gu{આપવું}} (\emph{dhyān āpvũ}), to give attention, is built from \emph{dhyān}, Sanskrit \emph{dhyāna}, \textquotedblleft{}meditation\textquotedblright{}, taken over as a tatsama (CDIAL 6889), and the inherited verb \emph{āpvũ}, \textquotedblleft{}to give\textquotedblright{}.
\end{cousinweb}

\begin{grammarlens}[title={What you've built}]
One more word for everyday life at A2.
\end{grammarlens}

\subsection*{Your turn}
\begin{itemize}
\raggedright
  \item \emph{Say it:} \emph{dhyān āpvũ}
  \item \emph{Say it:} \emph{dhyān āpvũ}, once more
  \item \emph{Say it:} \emph{vidyārthī vargmā̃ dhyān āpe che} — The students pay attention in class.
  \item \emph{Recall:} say the Gujarati for translation, then the Gujarati for dictionary, then say \emph{dhyān āpvũ} again
\end{itemize}

\begin{tcolorbox}[breakable,colback=teal!4,colframe=teal!35!black,title={Before you move on}]
What does \gu{ધ્યાન} \gu{આપવું} mean? (\textquotedblleft{}To pay attention\textquotedblright{}.) Say it once more.
\end{tcolorbox}
\section[\texorpdfstring{\gu{યાદ} \gu{કરવું} (yād karvũ)}{યાદ કરવું (yād karvũ)}]{\gu{યાદ} \gu{કરવું} (yād karvũ) — to memorize, to recall}

\label{lesson:GU-C251-yad-karvu}



\begin{quote}
Before the new one: say the Gujarati for dictionary, then the Gujarati for to pay attention.
\end{quote}

\subsection*{You'll want to know: \gu{યાદ} \gu{કરવું}}
\textbf{\gu{યાદ} \gu{કરવું}} — \emph{yād karvũ} — \textquotedblleft{}to memorize, to recall\textquotedblright{}.

\begin{quote}
\textbf{\gu{હું} \gu{રોજ} \gu{દસ} \gu{નવા} \gu{શબ્દ} \gu{યાદ} \gu{કરું} \gu{છું}.} — \emph{hũ roj das navā śabda yād karũ chũ} — \textquotedblleft{}I memorize ten new words every day.\textquotedblright{}
\end{quote}

\begin{cousinweb}
\textbf{\gu{યાદ} \gu{કરવું}} (\emph{yād karvũ}) is built from \emph{yād}, \textquotedblleft{}memory\textquotedblright{}, a loan usually traced to Persian \emph{yād}, and the verb \emph{karvũ}, \textquotedblleft{}to do\textquotedblright{}.
\end{cousinweb}

\begin{grammarlens}[title={What you've built}]
One more word for everyday life at A2.
\end{grammarlens}

\subsection*{Your turn}
\begin{itemize}
\raggedright
  \item \emph{Say it:} \emph{yād karvũ}
  \item \emph{Say it:} \emph{yād karvũ}, once more
  \item \emph{Say it:} \emph{hũ roj das navā śabda yād karũ chũ} — I memorize ten new words every day.
  \item \emph{Recall:} say the Gujarati for dictionary, then the Gujarati for to pay attention, then say \emph{yād karvũ} again
\end{itemize}

\begin{tcolorbox}[breakable,colback=teal!4,colframe=teal!35!black,title={Before you move on}]
What does \gu{યાદ} \gu{કરવું} mean? (\textquotedblleft{}To memorize, to recall\textquotedblright{}.) Say it once more.
\end{tcolorbox}
